\subsection{Graded normalized projective spectra over arbitrary rings}
\label{subsec:normalized-projective-spectra}

This construction begins with the absolute real seminorm spectrum of a
polynomial ring over an arbitrary commutative coefficient ring. Prescribing
an arbitrary coefficient seminorm is then an actual subspace condition.
The construction does not assume a field or a bounded coefficient seminorm.

\begin{definition}[Actual graded normalized spectrum]
\label{projective:general-normalized-spectrum}
\lean{EnvelopeProperty.Berkovich.Spectrum.IsPolynomialGraded,
  EnvelopeProperty.Berkovich.Spectrum.isPolynomialGraded_iff,
  EnvelopeProperty.Berkovich.Spectrum.isClosed_setOf_isPolynomialGraded,
  EnvelopeProperty.Berkovich.Spectrum.isPolynomialGraded_iff_isGraded,
  EnvelopeProperty.Berkovich.NormalizedProjectiveSpectrum,
  EnvelopeProperty.Berkovich.NormalizedProjectiveSpectrum.instTopologicalSpace,
  EnvelopeProperty.Berkovich.NormalizedProjectiveSpectrum.instT2Space,
  EnvelopeProperty.Berkovich.NormalizedProjectiveSpectrum.continuous_val,
  EnvelopeProperty.Berkovich.NormalizedProjectiveSpectrum.isGraded,
  EnvelopeProperty.Berkovich.NormalizedProjectiveSpectrum.sup_val_X,
  EnvelopeProperty.Berkovich.NormalizedProjectiveSpectrum.val_X_le_one,
  EnvelopeProperty.Berkovich.NormalizedProjectiveSpectrum.continuous_iff,
  EnvelopeProperty.Berkovich.NormalizedProjectiveSpectrum.exists_val_X_eq_one,
  EnvelopeProperty.Berkovich.NormalizedProjectiveSpectrum.chartOpen,
  EnvelopeProperty.Berkovich.NormalizedProjectiveSpectrum.isOpen_chartOpen,
  EnvelopeProperty.Berkovich.NormalizedProjectiveSpectrum.exists_mem_chartOpen,
  EnvelopeProperty.Berkovich.NormalizedProjectiveSpectrum.isClosed_carrier}
\leanok
\uses{spectrum:affine-seminorm,spectrum:topology,spectrum:graded-support}
For $P=R[X_0,\ldots,X_n]$, a real seminorm is polynomial-graded when its
value on every polynomial is the maximum of its values on all homogeneous
components. The normalized projective spectrum consists of such points
whose maximum coordinate value is one. Its topology is induced from the
actual absolute spectrum. Each coordinate-nonvanishing subset is open and
these opens cover, because the finite maximum is attained and equals one.
\end{definition}
\begin{proof}
The defining identities on all polynomials are closed evaluation identities.
Ultrametricity applied to the true homogeneous decomposition gives one
inequality automatically, so polynomial-gradedness equivalently means
that every component value is at most the original value. Normalization
is a finite continuous maximum identity. Mathlib's actual total-degree
decomposition identifies its abstract graded projection with the usual
homogeneous-component operator, so this condition also agrees with the
generic projection definition of gradedness. A coordinate value is nonzero on
an open evaluation locus, and a coordinate attains the normalized maximum.
\end{proof}

\begin{theorem}[Genuine arbitrary-ring Gauss coordinate charts]
\label{projective:general-normalized-charts}
\lean{EnvelopeProperty.Berkovich.NormalizedProjectiveSpectrum.coordinateChartMap,
  EnvelopeProperty.Berkovich.NormalizedProjectiveSpectrum.continuous_coordinateChartMap,
  EnvelopeProperty.Berkovich.NormalizedProjectiveSpectrum.coordinateChartMap_val_X,
  EnvelopeProperty.Berkovich.NormalizedProjectiveSpectrum.chartBound,
  EnvelopeProperty.Berkovich.NormalizedProjectiveSpectrum.one_le_chartBound,
  EnvelopeProperty.Berkovich.NormalizedProjectiveSpectrum.chartBound_pos,
  EnvelopeProperty.Berkovich.NormalizedProjectiveSpectrum.chartRadius,
  EnvelopeProperty.Berkovich.NormalizedProjectiveSpectrum.chartRadius_pos,
  EnvelopeProperty.Berkovich.NormalizedProjectiveSpectrum.continuous_chartBound,
  EnvelopeProperty.Berkovich.NormalizedProjectiveSpectrum.continuous_chartRadius,
  EnvelopeProperty.Berkovich.NormalizedProjectiveSpectrum.rawChartPoint,
  EnvelopeProperty.Berkovich.NormalizedProjectiveSpectrum.continuous_rawChartPoint,
  EnvelopeProperty.Berkovich.NormalizedProjectiveSpectrum.rawChartPoint_val_X,
  EnvelopeProperty.Berkovich.NormalizedProjectiveSpectrum.rawChartPoint_val_X_self,
  EnvelopeProperty.Berkovich.NormalizedProjectiveSpectrum.chartBound_eq_sup_dehom,
  EnvelopeProperty.Berkovich.NormalizedProjectiveSpectrum.rawChartPoint_normalized,
  EnvelopeProperty.Berkovich.NormalizedProjectiveSpectrum.chartHomogenize_support_subset,
  EnvelopeProperty.Berkovich.NormalizedProjectiveSpectrum.rawChartPoint_val_homogeneousComponent,
  EnvelopeProperty.Berkovich.NormalizedProjectiveSpectrum.rawChartPoint_isGraded,
  EnvelopeProperty.Berkovich.NormalizedProjectiveSpectrum.coordinateChartInversePoint,
  EnvelopeProperty.Berkovich.NormalizedProjectiveSpectrum.coordinateChartInv,
  EnvelopeProperty.Berkovich.NormalizedProjectiveSpectrum.coordinateChartInv_coe,
  EnvelopeProperty.Berkovich.NormalizedProjectiveSpectrum.continuous_coordinateChartInversePoint,
  EnvelopeProperty.Berkovich.NormalizedProjectiveSpectrum.continuous_coordinateChartInv,
  EnvelopeProperty.Berkovich.NormalizedProjectiveSpectrum.coordinateChartMap_val_dehomCoord,
  EnvelopeProperty.Berkovich.NormalizedProjectiveSpectrum.chartBound_coordinateChartMap,
  EnvelopeProperty.Berkovich.NormalizedProjectiveSpectrum.chartRadius_coordinateChartMap,
  EnvelopeProperty.Berkovich.NormalizedProjectiveSpectrum.coordinateChartMap_val_dehom_homogeneousComponent,
  EnvelopeProperty.Berkovich.NormalizedProjectiveSpectrum.rawChartPoint_coordinateChartMap_eq,
  EnvelopeProperty.Berkovich.NormalizedProjectiveSpectrum.coordinateChartInv_coordinateChartMap,
  EnvelopeProperty.Berkovich.NormalizedProjectiveSpectrum.coordinateChartMap_coordinateChartInv,
  EnvelopeProperty.Berkovich.NormalizedProjectiveSpectrum.coordinateChartHomeomorph,
  EnvelopeProperty.Berkovich.NormalizedProjectiveSpectrum.coordinateChartHomeomorph_apply,
  EnvelopeProperty.Berkovich.NormalizedProjectiveSpectrum.coordinateChartHomeomorph_symm_apply,
  EnvelopeProperty.Berkovich.NormalizedProjectiveSpectrum.coordinateChartMap_val_C,
  EnvelopeProperty.Berkovich.NormalizedProjectiveSpectrum.rawChartPoint_val_C}
\leanok
\uses{projective:general-normalized-spectrum,spectrum:principal-open,spectrum:gauss-point,spectrum:gauss-joint-continuity,charts:homogeneous-substitution,charts:localized-homogeneous-substitution}
Every coordinate-nonvanishing open of the actual normalized projective
spectrum is homeomorphic to the absolute spectrum of
$R[U_i:i\ne j]$. The forward map evaluates every polynomial after its real
ratio substitution $U_i=X_i/X_j$ in the coordinate localization. The
inverse is the genuine Gauss seminorm of radius
$\max(1,\max_i|U_i|)^{-1}$ pulled back by $X_j=T$, $X_i=U_iT$.
\end{theorem}
\begin{proof}
The actual principal localization extends the seminorm uniquely when its
coordinate is nonzero. Pullback through the ratio homomorphism gives the
forward seminorm, continuously, with the actual coordinate ratios.

The Gauss construction works over any commutative ring. Its coefficient
of $T^d$ after homogeneous substitution is the dehomogenization of the
true degree-$d$ component. This proves polynomial-gradedness on every
polynomial. The finite maximum of the substituted coordinate values is
exactly the chosen bound, so multiplication by its positive inverse
normalizes it. Joint Gauss continuity and continuity of that positive
radius prove continuity of the inverse.

Normalization recovers the radius from the forward ratios. The genuine
homogeneous scaling identity in the coordinate localization then recovers
every component, and gradedness recovers the entire original seminorm.
For the other inverse, localize the substitution also at $T$. Its full
composition with the ratio homomorphism is the coefficient inclusion
on every affine polynomial; uniqueness of the actual seminorm extension
and the Gauss coefficient formula recover that polynomial's original
value. Thus both inverse identities are whole-point identities, with no
field, domain or generator-value inference.
\end{proof}

\begin{theorem}[Arbitrary prescribed-base projective spectra]
\label{projective:general-relative-normalized-charts}
\lean{EnvelopeProperty.Berkovich.RelativeNormalizedProjectiveSpectrum,
  EnvelopeProperty.Berkovich.RelativeNormalizedProjectiveSpectrum.instTopologicalSpace,
  EnvelopeProperty.Berkovich.RelativeNormalizedProjectiveSpectrum.instT2Space,
  EnvelopeProperty.Berkovich.RelativeNormalizedProjectiveSpectrum.isClosed_coefficientLocus,
  EnvelopeProperty.Berkovich.RelativeNormalizedProjectiveSpectrum.isClosedEmbedding_val,
  EnvelopeProperty.Berkovich.RelativeNormalizedProjectiveSpectrum.chartOpen,
  EnvelopeProperty.Berkovich.RelativeNormalizedProjectiveSpectrum.isOpen_chartOpen,
  EnvelopeProperty.Berkovich.RelativeNormalizedProjectiveSpectrum.exists_mem_chartOpen,
  EnvelopeProperty.Berkovich.RelativeNormalizedProjectiveSpectrum.chartSubtypeHomeomorph,
  EnvelopeProperty.Berkovich.RelativeNormalizedProjectiveSpectrum.coordinateChartHomeomorph,
  EnvelopeProperty.Berkovich.RelativeNormalizedProjectiveSpectrum.coordinateChartHomeomorph_val_apply,
  EnvelopeProperty.Berkovich.RelativeNormalizedProjectiveSpectrum.coordinateChartHomeomorph_val_X,
  EnvelopeProperty.Berkovich.RelativeNormalizedProjectiveSpectrum.coordinateChartHomeomorph_symm_coe,
  EnvelopeProperty.Berkovich.RelativeNormalizedProjectiveSpectrum.coefficientHomeomorph,
  EnvelopeProperty.Berkovich.RelativeNormalizedProjectiveSpectrum.isClosed_carrier,
  EnvelopeProperty.Berkovich.RelativeNormalizedProjectiveSpectrum.isCompact_carrier,
  EnvelopeProperty.Berkovich.RelativeNormalizedProjectiveSpectrum.instCompactSpace,
  EnvelopeProperty.Berkovich.RelativeNormalizedProjectiveSpectrum.zeroCoordinates}
\leanok
\uses{projective:general-normalized-charts,spectrum:relative-generating-polydisc}
Prescribing any real seminorm $v$ on $R$ gives a genuine closed coefficient
subspace of the normalized projective spectrum. Its coordinate charts are
homeomorphic to $\operatorname{RelativeSpectrum}_v R[U_i:i\ne j]$, preserving
every coefficient and every affine polynomial. This prescribed-base
normalized space is compact and Hausdorff, without a global bound on $v$.
\end{theorem}
\begin{proof}
Every coefficient equality is a closed evaluation identity, so their
intersection is the actual closed prescribed subspace. The whole ratio
homomorphism and homogeneous Gauss substitution preserve all constant
polynomials; the absolute coordinate homeomorphisms therefore restrict
to the actual fixed-coefficient loci, giving the relative chart comparison.

The coordinate bounds $|X_i|\le1$ define a genuine compact generating
polydisc in the arbitrary prescribed-base finite-type spectrum. For each
polynomial its true bound uses its finite coefficient seminorm maximum,
so the values of $v$ need not all be at most one. Gradedness and the finite
normalization identity are closed; hence their intersection with this
polydisc is compact. Hausdorffness follows from the actual evaluation
embedding of the seminorm spectrum. An actual coefficient-spectrum point is obtained by pulling the given
base seminorm back through evaluation of all affine variables at zero.
Apply the genuine inverse chart to obtain a normalized point; this supplies
a real witness even when the coefficient seminorm is nontrivial. The empty
affine family at projective dimension zero is handled by the same maximum
bound of one. No native Proj or analytic sheaf identification is included
in this compactness assertion.
\end{proof}
